\section{Wstęp}
% Aplikacje: MuJoCo, Blender,  VistaKine
%https://www.blender.org
%https://github.com/google-deepmind/mujoco
%https://www.vistakine.com/playground
Współczesny rynek oprogramowania oferuje wiele narzędzi umożliwiających modelowanie trójwymiarowe, symulację zjawisk fizycznych oraz analizę ruchu obiektów.
Jedną z najpopularniejszych aplikacji implementującej symulowanie dynamiki ciał jest Blender\cite{blender}. 
Jest to bardzo rozbudowna, otwartoźródłowa aplikacja wykorzystywana głównie do modelowania 3D, natomiast daje możliwość symulowania fizyki tworzonych scen.


adsfkljasdfjk
Wytyczne:
Wstęp nie może składać się z podrozdziałów; powinien składać się z następujących akapitów:
 I akapit: w zależności od tematyki pracy: ogólny stan rynku. Przykład: jeżeli praca polega
na stworzeniu aplikacji, pierwszy akapit ma opisać stan rynku podobnych programów
(dostępnych w chwili obecnej). Na każdą aplikację należy poświęcić 4-5 zdań (podając
odnośniki do literatury). Uwaga: te właśnie aplikacje zostaną szczegółowo omówione
w następnym rozdziale. W akapicie powinno znaleźć się jak najwięcej odniesień do
literatury.
 II akapit: główne cele pracy, cele dodatkowe. Wyszczególnienie celów w omawianym
akapicie powinno przyjąć postać: "Głównym celem pracy jest opracowanie...". Jeżeli
występuje kilka celów, można je wypunktować.
 III akapit: Struktura pracy jest następująca: w Rozdziale 1 dokonano szczegółowej
charakterystyki porównawczej... Rozdział 2 opisuje... W Rozdziale 3 przedstawiono.
Rozdział ostatni podsumowuje pracę.
